\documentclass[11pt, headheight=30pt]{scrartcl}
\usepackage[white, nologo, hw]{darkWhite}
\usepackage{mathrsfs}

\newcommand{\BBB}{\mathscr{B}}
\DeclareMathOperator{\cov}{cov}
\DeclareMathOperator{\var}{var}

\title{18.675 Problem Set 5}
\begin{document}
% PROBLEMS SET 5 FOR 18.675 (FALL 2023)
% Problem 1. For x ∈ [0, 1), set τ (x) = 2x mod 1. Show that τ is a measure-preserving
% transformation of ([0, 1), B([0, 1)), dx), and that τ is ergodic. Identify the invariant function f corresponding
% to each integrable function f.
\section{}
\begin{prob*}
    For $x \in [0, 1)$, set $\tau(x) = 2x \mod 1$.
    Show that $\tau$ is a measure-preserving
    transformation of $([0, 1), \BBB([0, 1)), dx)$, and that
    $\tau$ is ergodic. Identify the invariant function
    $\overline{f}$ corresponding to each integrable function $f$.
\end{prob*}

Alright, let's recall the definitions.
\begin{itemize}
    \item A \vocab{measure-preserving transformation} $\tau$ satisfies
    \[
        \mu \circ \tau^{-1} = \mu
    \]
    over $\BBB$.
    \item An \vocab{invariant subset} $A \in \BBB$ satisfies $\tau^{-1}(A) = A$;
    these form a sigma algebra $\BBB_\tau$.
    \item If $\BBB_\tau$ contains only sets of measure $0$ or $1$,
    then $\tau$ is \vocab{ergodic}.
\end{itemize}

First off, for all $0\leq a < b < 1$,
% \[
%     \mu\left(\tau^{-1}([a,b))\right) = \mu\left(\left[\frac{a}{2}, \frac{b}{2}\right)\right) + \mu\left(\left[\frac{a+1}{2}, \frac{b+1}{2}\right)\right) = \frac{b-a}{2} + \frac{b-a}{2} = b-a = \mu([a,b)). 
% \]
\begin{align*}
    \mu\left(\tau^{-1}([a,b))\right) &= \mu\left(\left[\frac{a}{2}, \frac{b}{2}\right)\right) + \mu\left(\left[\frac{a+1}{2}, \frac{b+1}{2}\right)\right) \\
    &= \frac{b-a}{2} + \frac{b-a}{2}= b-a = \mu([a,b)).
\end{align*}
Recall that $\BBB([0,1))$ is generated by the intervals $[a,b)$,
thus $\tau$ is measure-preserving.

Next, consider any invariant subset $\tau^{-1}(A) = A$. Let
\[
    c_k = \frac{1}{2\pi}\int_0^1 \id_A e^{-2\pi i k x} dx,\qquad k\in \ZZ
\]
be the \vocab{Fourier series} of $\id_A$.
Then,
\begin{align*}
    c_k &= \frac{1}{2\pi}\int_0^1 \id_A e^{-2\pi i k x} dx\\
    &= \frac{1}{2\pi}\int_0^1 \left(\id_A e^{-2\pi i k x}\right)\circ \tau dx \\
    &= \frac{1}{2\pi}\int_0^1 \id_{\tau^{-1}(A)} e^{-2\pi i k \tau(x)} dx \\
    &= \frac{1}{2\pi}\int_0^1 \id_{A} e^{-2\pi i k 2x} dx \\
    &= c_{2k}.
\end{align*}
Thus, $c_k = c_{2k} = c_{4k} = \cdots$ for all $k \in \ZZ$.
By Parseval's identity, we know
\[
    \sum_{k\in \ZZ} |c_k|^2 = \frac{1}{(2\pi)^2}\int_0^1 |\id_A|^2 dx = \frac{1}{(2\pi)^2}\mu(A)
\]
Now, the RHS is supposed to be finite. Thus, all $c_k$ for $k\neq 0$
must be $0$. Thus, $\id_A$ must be constant almost everywhere (by considering $\id_A - 2\pi c_0$).
In other words, $A$ must have measure $0$ or $1$; thus, $\tau$ is ergodic.

Finally, we want to identify the invariant function
$\overline{f}$ corresponding to an integrable function $f$.
\idea{
    This is obviously
    just $\EE[f]$; I guess you want me to show that this actually
    works. Let's blow it out the water with the like, 17 theorems
    we know about this.
}

I interpret this to mean the almost-sure limit of $S_n / n$
guaranteed by Birkhoff's ergodic theorem, which by von
Neumann's $L^1$ ergodic theorem is also the $L^1$ limit
of $\frac{S_n}{n}$. I already know a priori that you can
take $\overline{f}$ to be constant, because invariant
functions on ergodic systems are constant almost everywhere.
So let $\overline{f} = c$ or something. Because this is the $L^1$
limit, I know that
\[
    \int_{[0,1)} \abs{\frac{S_n}{n} - c} \geq \abs{\int_{[0,1)} \frac{S_n}{n} - c} 
\]
but the LHS converges to $0$. Thus $c$ is the limit of
\[
    \int_{[0,1)} \frac{S_n}{n} = \frac{1}{n}\sum_{k=1}^n \int_{[0,1)} f\circ \tau^k = \frac{1}{n}\sum_{k=1}^n \int_{[0,1)} f = \frac{1}{n}\sum_{k=1}^n \EE[f] = \EE[f].    
\]
% lol you can't be serious rn.
% Problem 2. Fix a ∈ [0, 1) and define, for x ∈ [0, 1), τ (x) = x + a mod 1. Show that τ is
% also a measure-preserving transformation of ([0, 1), B([0, 1)), dx). Determine for which values of a
% the transformation τ is ergodic. Hint: you may use the fact that any integrable function f on [0, 1)
% whose Fourier coefficients all vanish must itself vanish a.e., where the Fourier coefficients (ck)k∈Z
% are defined by ck =
% 1
% 2π
% R 1
% 0
% f(x)e
% −2πikxdx for all k ∈ Z. Identify, for all values of a, the invariant
% function f corresponding to an integrable function f.
\section{}
\begin{prob*}
    Fix $a \in [0, 1)$ and define, for $x \in [0, 1)$, $\tau(x) = x + a \mod 1$.
    Show that $\tau$ is also a measure-preserving transformation of
    $([0, 1), \BBB([0, 1)), dx)$.
    Determine for which values of $a$ the transformation $\tau$ is ergodic.
    % Hint: you may use the fact that any integrable function $f$ on $[0, 1)$
    % whose Fourier coefficients all vanish must itself vanish a.e.,
    % where the Fourier coefficients $(c_k)_{k \in \ZZ}$ are defined by
    % $c_k = \frac{1}{2\pi} \int_0^1 f(x) e^{-2\pi i k x} dx$ for all $k \in \ZZ$.
    Identify, for all values of $a$, the invariant function $f$ corresponding
    to an integrable function $f$.
\end{prob*}

% Trivial.

$\tau$ is measure-preserving because uh
% \[
%     \mu(\tau^{-1}(E)) = \mu(\tau^{-1}(E \cap [0, 1-a))) + \mu(\tau^{-1}(E \cap [1-a, 1))) = \mu(E \cap [0, 1-a)) + \mu(E \cap [1-a, 1)) = \mu(E).
% \]
\begin{align*}
    \mu(\tau^{-1}(E)) &= \mu(\tau^{-1}(E \cap [0, 1-a))) + \mu(\tau^{-1}(E \cap [1-a, 1))) \\
    &= \mu(E \cap [0, 1-a)) + \mu(E \cap [1-a, 1)) = \mu(E).
\end{align*}
where the second equality is because the Lebesgue measure
is defined to be translation-invariant.

Obviously $\tau$ is ergodic iff $a$ is irrational; let me prove it.

\subsection{$a$ is rational}

If $\tau = p/q$ when reduced where $p,q\in \ZZ$, then the set
\[
    \left(\frac{m}{q}, \frac{m + 1/2}{q}\right),\qquad 0\leq m \leq q-1    
\]
is invariant, and has measure $1/2$. In this case, the invariant function
can be explicitly constructed:
\[
    \frac{S_{nq}}{nq}(x) = \frac{1}{nq}\sum_{k=1}^{nq} f\circ \tau^k(x) = \frac{1}{nq}\sum_{k=1}^{nq} f\left(x + \frac{kp}{q}\right)
     = \frac{1}{q}\sum_{k=0}^{q-1} f\left(x + \frac{k}{q}\right)
\]

\subsection{$a$ is irrational}

Now, pick some invariant $A$ and define the Fourier series of $\id_A$ as before. Then,
\begin{align*}
    c_k &= \frac{1}{2\pi}\int_0^1 \id_A e^{-2\pi i k x} dx\\
    &= \frac{1}{2\pi}\int_0^1 \left(\id_A e^{-2\pi i k x}\right)\circ \tau dx \\
    &= \frac{1}{2\pi}\int_0^1 \id_{\tau^{-1}(A)} e^{-2\pi i k \tau(x)} dx \\
    &= \frac{1}{2\pi}\int_0^1 \id_{A} e^{-2\pi i k (x+a)} dx \\
    &= \frac{1}{2\pi}\int_0^1 \id_{A} e^{-2\pi i k x} dx e^{-2\pi i k a} \\
    &= c_k e^{-2\pi i k a}.
\end{align*}
$a$ is irrational, thus  $e^{-2\pi i ka} \neq 1$ for all $k\neq 0$,
thus $c_k = 0$ for all $k \neq 0$, and then $f$ is
constant almost everywhere, thus $A$ has measure $0$ or $1$.
Thus $\tau$ is ergodic. 
% iff $a$ is irrational.

Finally, we want to identify the invariant function. % in the case where $a$ is irrational.
By identical reasoning to problem $1$, $\overline{f}$ must be constant almost everywhere,
and $\overline{f} = \EE[f]$.

% Problem 3. Call a sequence of random variables (Xn : n ∈ N) on a probability space stationary if
% for each n, k ∈ N the random vectors (X1, ..., Xn) and (Xk+1, ..., Xk+n) have the same distribution:
% for A1, ..., An ∈ B,
% \PP[X1 ∈ A1, ..., Xn ∈ An] = \PP[Xk+1 ∈ A1, ...., Xk+n ∈ An].
% Show that, if (Xn : n ∈ N) is a stationary sequence and X1 ∈ L
% p
% , for some p ∈ [1, ∞), then
% 1
% n
% Xn
% i=1
% Xi → X a.s. and in L
% p
% ,
% for some random variable X ∈ L
% p and find \EE[X].
\section{}
\begin{prob*}
    Call a sequence of random variables $(X_n)_{n \in \NN}$ on a probability space
    stationary if for each $n, k \in \NN$ the random vectors
    $(X_1, \ldots, X_n)$ and $(X_{k+1}, \ldots, X_{k+n})$ have the same distribution:
    for $A_1, \ldots, A_n \in \BBB$,
    \[
        \PP[X_1 \in A_1, \ldots, X_n \in A_n] = \PP[X_{k+1} \in A_1, \ldots, X_{k+n} \in A_n].
    \]
    Show that, if $(X_n)_{n \in \NN}$ is a stationary sequence and $X_1 \in L^p$,
    for some $p \in [1, \infty)$, then
    % $\frac{1}{n} \sum_{i=1}^n X_i \to X$ a.s. and in $L^p$,
    \[
        \frac{1}{n} \sum_{i=1}^n X_i \to X \text{ a.s. and in } L^p,
    \]
    for some random variable $X \in L^p$ and find $\EE[X]$.
\end{prob*}

Okay, I can push the measure forward by considering the map $X:\Omega \to \RR^\NN$
via $\omega \mapsto (X_1(\omega), X_2(\omega), \ldots)$; here, $\RR^\NN$
is equipped with the product sigma algebra (generated by the cylinder sets
with finitely many coordinates fixed); then let $\mu = \PP\circ X^{-1}$.

Consider the shift map $\theta: \Omega \to \Omega$ defined by
$\theta(\omega) = (\omega_2, \omega_3, \ldots)$.
Then, the condition precisely states that $\theta$ is measure-preserving;
it preserves the measure on all cylinder sets, and these generate the
product sigma algebra.

Thus, von Neumann's $L^p$ ergodic theorem applied to the coordinate function
$X_1\circ X^{-1}: \RR^{\NN} \to \RR$ shows $\frac{S_n}{n}\to X$ AS and in $L^p$,
where $S_n = X_1 + \cdots + X_n$, for some rv $X$.

Note that $L^p\subset L^1$ for all $p\in [1,\infty)$ so $X_1 \in L^1$ as well,
and $\frac{S_n}{n}\to X$ in $L^1$ as well.

Observe that
\[
    \EE\left[\frac{X_1 + \dots + X_n}{n}\right] = \frac{1}{n}\sum_{i=1}^n \EE[X_i] = \EE[X_1].    
\]
for all $n$. Also,
\[
    \abs{\EE[X] - \EE\left[\frac{S_n}{n}\right]} \leq \int_{\Omega} \abs{X - \frac{S_n}{n}} d\mu \to 0.    
\]
Thus, $\EE[X] = \EE[X_1]$.



% Problem 4. Let (Xn : n ∈ N) be a sequence of independent random variables, such that
% \EE[Xn] = µ and \EE[X4
% n
% ] ≤ M for all n, for some constants µ ∈ R and M < ∞. Set Pn =
% X1X2 + X2X3 + · · · + Xn−1Xn. Show that Pn/n converges a.s. as n → ∞ and identify the limit.
\section{}
\begin{prob*}
    Let $(X_n)_{n \in \NN}$ be a sequence of independent random variables, such that
    $\EE[X_n] = \mu$ and $\EE[X_n^4] \leq M$ for all $n$, for some constants $\mu \in \RR$ and $M < \infty$.
    Set $P_n = X_1 X_2 + X_2 X_3 + \cdots + X_{n-1} X_n$.
    Show that $P_n/n$ converges a.s. as $n \to \infty$ and identify the limit.
\end{prob*}

We claim $\frac{P_n}{n}\to \mu^2$ almost surely.

The trick is to split the sequence into even and odd components;
\[
    P_n = (X_1X_2 + X_3X_4 + \dots) + (X_2X_3 + X_4X_5 + \dots)    
\]
Each of these subsequences contains independent terms. The mean
of each term is $\EE[X_iX_{i + 1}] = \EE[X_i]\EE[X_{i+1}] = \mu^2$,
and the fourth moment is bounded by $M^2$ by independence.

Thus, theorem $10.1.1$ says
\[
    \frac{X_1X_2 + \dots + X_{n-1}{X_n}}{n/2}\to \mu^2    
\]
almost surely; we are ignoring off-by-one errors due to the parity
of $n$ because those do not matter in the limit. Thus,
the average of the entire sum converges a.s. to $\mu^2$ as well.

% First, let $Y_n = X_n - \mu$; then
% $\EE[Y_n^4]\leq 2^4(\EE[X_n^4] + \mu^4)  < \infty$, and
% \[
%     P_n = Y_1Y_2 + \dots + Y_{n-1}Y_n + \underbrace{\mu(X_1 + \dots + X_{n-1}) + \mu(X_2 + \dots + X_n)}_\star - \mu^2(n-1).
% \]
% Well, $\star/n$  goes to $2\mu^2$ almost surely by SLLN, and $\frac{\mu^2(n-1)}{n} \to \mu^2$ almost surely by brain cells,
% thus it suffices to show that $\frac{Y_1Y_2 + \dots + Y_{n-1}Y_n}{n} \to 0$ almost surely. So we reduce to the $0$-mean case.
% Redefine $P_n = Y_1Y_2 + \dots + Y_{n-1}Y_n$.

% We know $X_n, X^2_n, X^3_n, X^4_n$ are all integrable,
% and by independence, $\EE[X_iX_j^3] = \EE[X_iX_jX_k^2] = \EE[X_iX_jX_kX_l] = 0$ for all $i,j,k,l$.

% Now we consider
% \[
%     P_n^2 = (Y_1Y_2 + \dots + Y_{n-1}Y_n)^2 = \sum_{i=1}^{n-1} Y_i^2Y_{i+1}^2 + \sum_{i\neq j} Y_iY_jY_{i+1}Y_{j+1}. 
% \]
% The second sum has $0$ expectation, thus
% \[
%     \EE[P_n^2] = \sum_{i=1}^{n-1} \EE[Y_i^2Y_{i+1}^2] \leq \sum_{i=1}^{n-1} \EE[Y_i^4]^{1/2}\EE[Y_{i+1}^4]^{1/2} = M(n-1). 
% \]
% where the inequality is by Cauchy-Scwharz. Thus
% \[
%     \EE \left[\left(\frac{P_n}{n}\right)^2\right]\leq \frac{M(n-1)}{n^2}    
% \]
% \idea{
%     Considering the summation
%     \[
%         \EE \sum_n \left(\frac{P_n}{n}\right)^2 \leq \sum_n \frac{M(n-1)}{n^2} = \infty    
%     \]
%     we are sad because this doesn't tell us anything.
% }
% Instead, give me some rational $\delta > 0$. I wish to show that
% \[
%     \PP\left(\exists N : \forall m > N, \abs{\frac{P_m}{m}} > \delta\right) = 0.
% \]

% Instead, consider
% \[
%     \EE \sum_n \left(\frac{P_{a_n}}{a_n}\right)^2 \leq \sum_n \frac{M(a_n-1)}{a_n^2}   
% \]
% where $a_n = \floor{n^\alpha}$ and $\alpha > 1$ is some constant. Then, the RHS converges,
% and we discover that
% \[
%     \sum_n \left(\frac{P_{a_n}}{a_n}\right)^2 < \infty \text{ almost surely.}    
% \]
% i.e. $P_{a_n}/a_n \to 0$ almost surely. Now, consider all rational numbers $\alpha > 1$;
% there are countably many of these, thus, it is almost surely true that
% for \emph{all} rational numbers $\alpha > 1$, $P_{a_n}/a_n \to 0$.


% Problem 5. The Cauchy distribution has density function f(x) = π
% −1
% (1 + x
% 2
% )
% −1
% for x ∈ R.
% Show that the corresponding characteristic function is given by ϕ(u) = e
% −|u|
% . Show also that, if
% X1, ..., Xn are independent Cauchy random variables, then the random variable (X1 + · · · + Xn)/n
% is also Cauchy.
\section{}
\begin{prob}
    The Cauchy distribution has density function $f(x) = \frac{1}{\pi} \frac{1}{1 + x^2}$ for $x \in \RR$.
    Show that the corresponding characteristic function is given by $\phi(u) = e^{-|u|}$.
\end{prob}

Okay, let us compute
\[
    \phi(u) = \EE[e^{iuX}] = \int_{-\infty}^\infty e^{iux} \frac{1}{\pi} \frac{1}{1 + x^2} dx.
\]
This is clearly symmetric about $u\to -u$ by sending $x\to -x$. Thus we
consider the case $u \geq 0$ only.

We can compute this by contour integration. Consider the contour $\gamma$
composed of the semi-circle $\gamma_R$ of radius $R$ in the upper half-plane
centered at $0$, and the line segment $[-R, R]$ on the real axis. In the limit,
the integral over the semi-circle is bounded by
\begin{align*}
    \abs{\int_{\gamma_R} e^{iuz} \frac{1}{\pi} \frac{1}{1 + z^2} dz}
    &\leq \pi R \cdot \sup_{z\in \gamma_R} \abs{e^{iuz}\frac{1}{\pi} \frac{1}{1 + z^2}}\\
    &\leq \pi R \cdot \sup_{z\in \gamma_R} \underbrace{\abs{e^{-u \Im(z)}}_{ < 1}\frac{1}{\pi} \frac{1}{1 + z^2}}\\
    &\leq \pi R \cdot \frac{1}{\pi} \frac{1}{R^2 - 1}\\
    &\to 0
\end{align*}
thus the whole contour integral becomes the integral over the real line.

On the other hand, by the residue theorem, the integral over the contour
is $2\pi i$ times the sum of the residues of the poles inside the contour.
The poles are at $z = \pm i$, and the residue at $z = i$ is
\[
    \lim_{z\to i} (z - i) e^{iuz} \frac{1}{\pi} \frac{1}{1 + z^2} = \frac{e^{-u}}{2i}
\]
Thus, the integral is $2\pi i \cdot \frac{e^{-u}}{2i} = \pi e^{-u}$.
Thus, $\phi(u) = e^{-|u|}$, as desired.

\begin{prob}
    Show also that, if $X_1, \ldots, X_n$ are independent Cauchy random variables,
    then the random variable $Y = (X_1 + \cdots + X_n)/n$ is also Cauchy.
\end{prob}

We compute the characteristic function of $Y$:
\[
    \phi_Y(u) = \EE[e^{iuY}] = \EE[e^{iu(X_1 + \cdots + X_n)/n}] = \EE[e^{iuX_1/n}]^n = \phi(u/n)^n = e^{-|u|}.
\]
Thus, $Y$ is Cauchy.

% Problem 6. For each n ∈ N, there is a unique probability measure on the sphere S
% n−1 =
% {x ∈ R
% n
% : |x| = 1} such that µn(A) = µn(UA) for all Borel sets A and all orthogonal n × n matrices
% U. Fix k ∈ N and, for n ≥ k, let γn denote the probability measure on R
% k which is the law of
% √
% n(x
% 1
% , ...., xk
% ) under µn. Show that
% (i) if X is a mean-zero Gaussian with covariance matrix given by In then X/|X| has law given by
% µn.
% (ii) if (
% p
% Xn : n ∈ N) is a sequence of independent N(0, 1) random variables and if Rn =
% X2
% 1 + · · · + X2
% n
% then Rn/
% √
% n → 1 a.s.
% (iii) γn converges weakly to the standard Gaussian distribution on R
% k as n → ∞.
% 1
\section{}
\begin{prob}
    For each $n \in \NN$, there is a unique probability measure on the sphere $S^{n-1} = \{x \in \RR^n : |x| = 1\}$
    such that $\mu_n(A) = \mu_n(UA)$ for all Borel sets $A$ and all orthogonal $n \times n$ matrices $U$.
    Fix $k \in \NN$ and, for $n \geq k$, let $\gamma_n$ denote the probability measure on $\RR^k$
    which is the law of $\sqrt{n} (x_1, \ldots, x_k)$ under $\mu_n$.
    
    Show that if $X$ is a mean-zero Gaussian with covariance matrix given by $I_n$ then $X/|X|$ has law given by $\mu_n$.
    % \begin{enumerate}[(i)]
    %     \item 
    %     \item if $(X_n)_{n \in \NN}$ is a sequence of independent $N(0, 1)$ random variables and if $R_n = \sqrt{X_1^2 + \cdots + X_n^2}$
    %         then $R_n/\sqrt{n} \to 1$ a.s.
    %     \item $\gamma_n$ converges weakly to the standard Gaussian distribution on $\RR^k$ as $n \to \infty$.
    % \end{enumerate}
\end{prob}

Observe that the density function (and law) of $X$ is in fact symmetric
under orthogonal transformations;
\[
    \mathcal{N}(0, I_n)(Ux) = \frac{1}{(2\pi)^{n/2}} e^{-\frac{1}{2} (Ux)^T Ux} = \frac{1}{(2\pi)^{n/2}} e^{-\frac{1}{2} x^T x} = \mathcal{N}(0, I_n)(x).
\]
Thus, the law of $X/|X|$ is also symmetric under orthogonal transformations,
and is a probability measure on the sphere $S^{n-1}$. Thus, it must be $\mu_n$.

\begin{prob}
    Show that if $(X_n)_{n \in \NN}$ is a sequence of independent $N(0, 1)$ random variables and if $R_n = \sqrt{X_1^2 + \cdots + X_n^2}$
    then $R_n/\sqrt{n} \to 1$ a.s.
\end{prob}

The sequence $(X^2_n)_{n \in \NN}$ is a sequence of i.i.d random variables
with mean $1$ (this is just the $\chi^2$ distribution).

Thus, by the strong law of large numbers, $\frac{X_1^2 + \cdots + X_n^2}{n} \to 1$
almost surely. Taking the square root of each side,
$\frac{R_n}{\sqrt{n}} \to 1$ almost surely.

\begin{prob}
    Show that $\gamma_n$ converges weakly to the standard
    Gaussian distribution on $\RR^k$ as $n \to \infty$.
\end{prob}
By part $1$,
\[
    \gamma_n \sim \sqrt{n} \frac{X_k}{\abs{X_n}}
\]
where $X_i = (x_1,\dots x_i)$ and each $(x_i)_{i\geq n}$ is a standard Gaussian.

By the previous problem, it is also true that
\[
    \frac{\abs{X_n}}{\sqrt{n}} \to 1 \text{ almost surely.}
\]
so it also converges in probability. Thus, by Slutsky's theorem,
$\gamma_n \to \mathcal{N}(0, I_k)$ in distribution.

% It suffices to show that the characteristic functions converge pointwise.

% 2 PROBLEMS SET 5 FOR 18.675 (FALL 2023)
% Problem 7. Let X and Y be integrable random variables and suppose that
% \EE[X | Y ] = X and \EE[Y | X] = X a.s.
% Show that X = Y a.s.
% Hint: Consider quantities like \EE[(X − Y )1X>c,Y ≤c] + \EE[(X − Y )1X≤c,Y ≤c].
\section{}
\setcounter{prob*}{6}
\begin{prob*}
    Let $X$ and $Y$ be integrable random variables and suppose that
    $\EE[X | Y ] = Y$ and $\EE[Y | X] = X$ a.s.
    Show that $X = Y$ a.s.
    Hint: Consider quantities like $\EE[(X - Y) \id_{X > c, Y \leq c}] + \EE[(X - Y) \id_{X \leq c, Y \leq c}]$.
\end{prob*}

Well, by definition of expected value, $\EE[Y \id_{Y\in A}] = \EE[X \id_{Y\in A}]$
for all $A$ and $\EE[X \id_{X\in A}] = \EE[Y \id_{X\in A}]$ for all $A$. For all $c\in \RR$,
\begin{align*}
    \EE[(X - Y)\id_{X\geq c}\id_{Y\geq c}] &= \EE[(X-Y)\id_{X\geq c}] - \underbrace{\EE[(X-Y)\id_{X\geq c}\id_{Y < c}]}_{\geq 0} \\
    & \leq 0
\end{align*}
By analogous arguments, $\EE[(Y - X)\id_{X\geq c}\id_{Y\geq c}] \leq 0$. Thus
$\EE[(X-Y)\id_{X\geq c}\id_{Y\geq c}] = 0$. Running it back, this implies
$\EE[(X - Y)\id_{X\geq c}\id_{Y < c}] = 0$ for all $c$. I claim that this implies
$\PP(X \geq c \text{ and } Y < c) = 0$ for all $c$. Indeed, suppose this has positive
measure; then it is the countable union of $\PP(X \geq c \text{and} Y < c - \frac{1}{n})$,
thus one of these has positive measure; but then
\[
    \EE[(X - Y)\id_{X\geq c}\id_{Y < c - \frac{1}{n}}] \geq \frac1n \PP(X \geq c \text{ and } Y < c - \frac{1}{n}) > 0   
\]
contradiction. Okay! Thus $\PP(X \geq c \text{ and } Y < c) = 0$ for all $c$. Taking
the union across all rational $c$, $\PP(X > Y) = 0$. By symmetry, $\PP(X < Y) = 0$.
Thus $X = Y$ almost surely.

% Problem 8. Let X, Y be two independent Bernoulli random variables with parameter p ∈ (0, 1).
% Let Z = 1X+Y =0. Compute \EE[X |Z] and \EE[Y |Z].
\section{}
\begin{prob*}
    Let $X, Y$ be two independent Bernoulli random variables with parameter $p \in (0, 1)$.
    Let $Z = \id_{X+Y = 0}$.
    Compute $\EE[X |Z]$ and $\EE[Y |Z]$.
\end{prob*}
You can just compute everything. $\PP(X = 0, Y = 0) = (1-p)^2$,
$\PP(X = 1, Y = 0) = p(1-p)$, $\PP(X = 0, Y = 1) = p(1-p)$,
and $\PP(X = 1, Y = 1) = p^2$. Thus,
\[
    \EE[X | Z] = \begin{cases}
        \frac{1}{2 - p} &\text{ if } Z = 0\\
        0 &\text{ if } Z = 1
    \end{cases}
\]
and ditto for $Y$.

% Problem 9. Give an example of a random variable X and two σ algebras H and G such that X is
% independent of H and G is independent of H but nevertheless
% E[X | σ(H, G)] 6= E[X | G].
\section{}
\begin{prob*}
    Give an example of a random variable $X$ and two $\sigma$-algebras $\mathfrak{H}$ and $\mathfrak{G}$ such that $X$ is
    independent of $\mathfrak{H}$ and $\mathfrak{G}$ is independent of $\mathfrak{H}$ but nevertheless
    $\EE[X | \sigma(\mathfrak{H}, \mathfrak{G})] \neq \EE[X | \mathfrak{G}]$.
\end{prob*}

Let $Y, Z$ be two independent fair coin flips and let $\mathfrak{G}$ and $\mathfrak{H}$ be
their sigma algebras. Then, $X = \id_{Y + Z = 1}$ is independent of $\mathfrak{G}$ and $\mathfrak{H}$,
but $\EE[X | \sigma(\mathfrak{G}, \mathfrak{H})] = X$ while $\EE[X | \mathfrak{G}] = \frac{1}{2}$.

\end{document}